\subsection{\colorbox{yellow!80}{\parbox{\textwidth}{PT06-CSP: Wymuszenie zmiany polityki CSP}}}
\subsubsection{Opis}
Atak miał na celu zmianę polityki CSP obowiązującej na stronie profilu użytkownika (/profile). Modyfikacja odbyła się poprzez odpowiednią preparację linku do awatara użytkownika.

Nie udało się znaleźć innych modułów aplikacji, gdzie prezentowane były awatary użytkownika, dlatego atak ten jest bardzo ograniczony.

\subsubsection{Warunki wykonania ataku}
Atakujący posiada poprawny token użytkownika.

\subsubsection{Szczegóły wykonania}
Żądanie wykonujące atak zostało przedstawione w tabeli.
\begin{table}[H]
\begin{tabular}{|l|l|l|l|l|l|}
\hline
\rowcolor[HTML]{EFEFEF} 
\multicolumn{1}{|c|}{\cellcolor[HTML]{EFEFEF}\textbf{Metoda}} & \multicolumn{1}{c|}{\cellcolor[HTML]{EFEFEF}\textbf{Końcówka}} & \multicolumn{1}{c|}{\cellcolor[HTML]{EFEFEF}\textbf{Nagłowki}} & \multicolumn{1}{c|}{\cellcolor[HTML]{EFEFEF}\textbf{Ciasteczka}} & \multicolumn{1}{c|}{\cellcolor[HTML]{EFEFEF}\textbf{Dane}}                             & \multicolumn{1}{c|}{\cellcolor[HTML]{EFEFEF}\textbf{Odpowiedź/wynik}} \\ \hline
POST                                                          & /profile/image/url                                        & \begin{tabular}[c]{@{}l@{}}Authentication \\ Bearer TOKEN \end{tabular}                                   & \begin{tabular}[c]{@{}l@{}}token: \\ TOKEN\end{tabular}                    & \begin{tabular}[c]{@{}l@{}}\{"imageUrl":"http://assets/public/images/uploads/1.jpg;\\ script-src 'unsafe-inline'"\}\end{tabular} & 302 /profile                               \\ \hline
\end{tabular}
\end{table}

\subsubsection{Skutki}

Po wykonaniu zapytania nagłówek CSP zwracany przez sklep na profilu użytkownika wyglądał następująco:
\begin{lstlisting}
img-src 'self' http://assets/public/images/uploads/1.jpg; script-src 'unsafe-inline'; script-src 'self' 'unsafe-eval' https://code.getmdl.io http://ajax.googleapis.com
\end{lstlisting}
co wskazuje na to, że \textbf{skrypty o niezweryfikowanym źródle zostały aktywowane} (co potwierdza obecność 'unsafe-inline' na liście script-src) na stronie profilu. 

\subsubsection{Zalecenia}
W celu zapobiegnięcia modyfikacji CSP należało by zaimplementować poprawioną metodę sprawdzania, czy podana jako link do zdjęcia wartość faktycznie jest linkiem.
Jedną z metod weryfikacji URL w JS jest tak samo nazwana biblioteka, która zwraca błąd jeśli tekst nie jest poprawnym łączem. 

Ponadto, ciekawym dodatkowym zabezpieczeniem jest usunięcie wszystkiego po znaku średnika. Teoretycznie może on występować w URI jednak jest on w pewien sposób "specjalny" i ma inne znaczenie niż zwykłe litery i cyfry (\href{https://www.rfc-editor.org/rfc/rfc3986\#section-2.2}{https://www.rfc-editor.org/rfc/rfc3986\#section-2.2}), dlatego i tak nie powinien być obecny zbyt często w łączach do obrazków. Takie zabezpieczenie automatycznie usunie wszystko co mogło by ingerować w politykę w CSP, więc nawet jeśli weryfikacja URL nie zadziała to do CSP nie zostanie wstawione nic co umożliwi XSS.

Przykład poprawki modyfikacji CSP w pliku profileImageUpload.ts:
\begin{lstlisting}
const { URL } = require('url');
[...]
if (loggedInUser) {
    var myurl = ""
    try {
        myurl = new URL(requrl).toString();
        myurl = myurl.split(';')[0] 
    } catch (error) {
        myurl = "";
        next(error)
    }
[...]
\end{lstlisting}
Poprawka częściowo zainspirowana przez wpis \href{https://stackoverflow.com/questions/30931079/validating-a-url-in-node-js}{https://stackoverflow.com/questions/30931079/validating-a-url-in-node-js}.


\subsubsection{StoredXSS}
\subsection{\colorbox{green!80}{\parbox{\textwidth}{PT07-XSS: Prosty atak Stored XSS}}}
\subsubsection{Opis}

Atak polegał na celowej preparacji danych (próby objęły głównie pseudonim użytkownika) w taki sposób, aby powodował on wykonanie dowolnego skryptu po jego wyświetleniu na witrynie.

W badanej implementacji sanityzacja została poprawiona, dlatego nie udało się wykonać prostego XSS na pseudonimie użytkownika. Sprawdzone zostało wiele kombinacji, jednak ze względu na brak powodzenia można uznać, że użyta biblioteka jest na tyle dobra, że walka z nią nie ma większego sensu i lepiej spróbować podejść do problemu w innym miejscu, bądź inną, mniej bezpośrednią metodą.

W raporcie zostały opisane podjęte próby.

\subsubsection{Warunki wykonania ataku}
Różne w zależności od wersji ataku. Dla ataku na pseudonim atakujący musiał posiadać poprawny token użytkownika i wykonać atak na CSP pozwalający na uruchomienie dodatkowych skryptów.

\subsubsection{Szczegóły wykonania (podjęte próby)}

W tabeli zawarty został schemat prostego ataku XSS na pseudonim użytkownika z przykładowymi danymi. Badane były również takie wartości pola username jak np. "<<sscript>sscript>alert('Modified CSP')</script>", które działało w podstawowej implementacji JuiceShop, jednak tutaj zostało poprawnie zsanityzowane.

\begin{table}[H]
\begin{tabular}{|l|l|l|l|l|l|}
\hline
\rowcolor[HTML]{EFEFEF} 
\multicolumn{1}{|c|}{\cellcolor[HTML]{EFEFEF}\textbf{Metoda}} & \multicolumn{1}{c|}{\cellcolor[HTML]{EFEFEF}\textbf{Końcówka}} & \multicolumn{1}{c|}{\cellcolor[HTML]{EFEFEF}\textbf{Nagłowki}} & \multicolumn{1}{c|}{\cellcolor[HTML]{EFEFEF}\textbf{Ciasteczka}} & \multicolumn{1}{c|}{\cellcolor[HTML]{EFEFEF}\textbf{Dane}}                             & \multicolumn{1}{c|}{\cellcolor[HTML]{EFEFEF}\textbf{Odpowiedź/wynik}} \\ \hline
POST                                                          & /profile                                        & \begin{tabular}[c]{@{}l@{}}Authentication \\ Bearer TOKEN \end{tabular}                                   & \begin{tabular}[c]{@{}l@{}}token: \\ TOKEN\end{tabular}                    & \begin{tabular}[c]{@{}l@{}}\{'username': "\></p\><img\><a 
\\    
href=javascript:'javascript:alert('ModifiedCSP')'\></img\>"\}\end{tabular} & 302 /profile                               \\ \hline
\end{tabular}
\end{table}

Poszukiwania podatności na Stored XSS objęły również inne endpointy. Na każdym z nich testowane były różne kombinacje znaków (w szczególności /,<,>,",'), tagów script/img i innych losowych ciągów z jednoczesną obserwacją, czy wstawiane do bazy elementy wykazują potencjał na skuteczny XSS. W skład prób wchodziły:
\begin{itemize}
    \item Modyfikacja pola email przy rejestracji (POST /api/Users/):
    
    \lstinline[]|"email":"x<script>alert()</script>de@mail.com")|
    \item Dodawanie i modyfikacja pól message oraz author recenzji (PUT /rest/products/1/reviews: 
    
    \lstinline[]|"message":"<script<<<<script>alert()</script>"|,
    
    \lstinline[]|"author":"<<<script<<<script>alert()</script>"|)
    \item Dodawanie opini customer feedback z atakiem w polu comment
    \item Atak inspirowany przez mutation XSS na username (POST /profile:

    \lstinline[]|username=script<style = width:0px>><<|
     
    wydawał się częściowo skuteczne, ale znak < był zawsze usuwany, więc ostatecznie nic nie udało się osiągnąć 
\end{itemize}

\subsubsection{Skutki}
Na tym etapie \textbf{nie udało się wykonać skutecznego ataku XSS}. 


\subsection{\colorbox{green!80}{\parbox{\textwidth}{PT08-SQL-XSS: Atak Stored XSS wspierany przez SQL Injection}}}
\subsubsection{Opis}
Po licznych nieudanych próbach podstawowego ataku XSS uznano, że  jest on zbyt trudny w realizacji ze względu na poprawioną sanityzację. Jako formę jej ominięcia podjęto próbę modyfikacji zawartości bazy danych poprzez SQL Injection tak, aby uzyskać możliwość wstawienia ataku XSS bezpośrednio do niej (poprzez użycie SELECT CHAR(60) itp.). 
\subsubsection{Warunki wykonania ataku}
W zależności od podejmowanych prób w niektórych był wymagany poprawny token użytkownika.
\subsubsection{Szczegóły wykonania (próby)}
Podjęte zostały następujące próby realizacji ataku:
\begin{itemize}
    \item Wstrzyknięcie znaku bezpośrednio do bazy przez SQL Injection (to co opisany wcześniej SQL injection  GET /rest/products/search: 
    
    \begin{lstlisting}
        ?q=apple')); update users set username="xs" where id=23; --")
    \end{lstlisting}
    W trakcie prób do tego punktu udało się ustalić, że sanityzacja XSS dla pseudonimu odbywa się przed wstawieniem do bazy (nie było jednak wiadomo czy tylko przed, czy zarówno przed jak i po)
    \item zapytania aktualizujące i wstawiające opinie (PUT /rest/products/1/reviews)
    \item żądania zmiany pseudonimu (POST /profile)
\end{itemize}

\subsubsection{Skutki}
Podobnie jak w przypadku podstawowego XSS \textbf{nic nie udało się osiągnąć}.
